\section{Introduction}\label{ch2:sec:introduction}


Government announcements materially influence interest-rate markets (\cite{vergote2012}; \cite{rosa2013financial}; \cite{jubinski2013fomc}). Among these, the Federal Open Market Committee (FOMC) plays a particularly critical role in shaping both the U.S. and global economic outlook through its monetary policy decisions—most notably, adjustments to the target range of the federal funds rate, which have direct and immediate effects on the government bond market. Moreover, U.S. government bond derivatives, such as the 10-Year Treasury Note futures and options, are traded nearly continuously, with only a one-hour pause from 4 p.m. to 5 p.m. Central Time (CT). This continuous trading structure enables market participants to incorporate policy announcements and macroeconomic information into prices in real time, thereby making these instruments an ideal channel through which to observe the market’s instantaneous response to government communications.


Meanwhile, current models for predicting bond prices, grounded in the Efficient Market Hypothesis (EMH, \cite{fama1970}), rely primarily on historical quantitative data. For example, parametric models like the Heston Stochastic Volatility Model (\cite{Heston:1993}) and the Stochastic Alpha Beta Rho (SABR) Model (\cite{sabr2002}) must be calibrated to the current shape of the implied volatility surface. Although these traditional forecasting models employ rigorous econometric techniques and quantitative indicators, they often struggle to incorporate unstructured textual information, such as FOMC meeting minutes and economic briefings, which frequently contain subtle but critical signals about future policy direction.



As a result, there is a pressing need for forecasting models that can effectively capture and interpret qualitative information, especially the textual sentiment driven market movements (\cite{bollen2011twitter}, \cite{ding2015deep},\cite{ding2016knowledge}, and \cite{xu2018stock}).


However, unlike quantitative data, qualitative data are inherently discontinuous, making historical data insufficient for analyzing a wide range of potential scenarios. For instance, the Federal Reserve has adjusted the federal funds rate only 130 times since the 1990s. To address this challenge, we propose leveraging a fine-tuned large language model (LLM) to generate synthetic texts across a variety of scenarios, thereby enhancing the capability of volatility models to incorporate qualitative inputs.




On the other hand, recent advances in neural networks and LLMs have led to remarkable performance improvements across a broad range of natural language processing (NLP) tasks. LLMs are pre-trained on vast and diverse text corpora, which can be adapted to domain-specific tasks and have already been used for synthetic data generation in multiple settings (\cite{schick2021generating}; \cite{peng2023instruction}; \cite{yu2024large}; \cite{li2021synthetic}; \cite{meoni2024generating}). In finance, synthetic datasets are also used for robustness testing (\cite{cao2022synthetic}; \cite{carvajal2022synthetic}).


Moreover, synthetic data have been widely accepted by financial researchers as powerful tools for robustness testing. For instance, \cite{cao2022synthetic} created a synthetic limit order book to evaluate forecasting models, while \cite{carvajal2022synthetic} used synthetic stock prices to train trading algorithms. Beyond time-series data, the financial industry manages vast volumes of structured documents, such as financial reports, Bank of England Monetary Policy Committee (MPC) meeting minutes, and Federal Reserve FOMC meeting minutes. The simulation capabilities of LLMs make them ideal for generating synthetic structured textual data in these contexts.

% [Codex 修改版] 2026-04-17: 在引言中提前交代证据边界，避免与后文“sample-level / simple baselines only”的口径脱节。
This chapter proposes a fine-tuned LLM framework for simulating monetary-policy communication and decision-relevant analysis. Conditioned on macroeconomic and financial indicators, the model is trained to generate (i) domain-consistent analytical narratives and (ii) institutionally styled text aligned with FOMC Minutes sections. We then evaluate whether the synthetic text is stylistically faithful, input-grounded, and economically meaningful through leave-one-out masking tests, sentiment--return regressions, and policy-decision prediction against historical and market-implied baselines. All canonical Chapter 2 experiments are organized around meeting-level fixed random 80/10/10 splits with a shared seed, so the reported evidence should be interpreted as held-out meeting validation rather than as a time-series extrapolation exercise.



\subsection*{Research Motivation}

The motivation for this research arises from the critical importance of accurately forecasting monetary policy decisions and the well-recognized limitations of traditional econometric models in capturing qualitative and contextual signals. Firstly, monetary-policy forecasting remains important for asset pricing and risk management. Secondly, conventional quantitative models do not fully exploit textual policy signals from minutes, statements, and staff discussions. As a result, they may miss informative cues about policy direction.


LLMs provide a practical route to close this gap by integrating structured indicators with unstructured text. Furthermore, evidence on using LLMs to simulate committee-style policy communication and decision-relevant reasoning is still limited. Hence, this chapter addresses that gap through a scenario-based simulation framework. More specifically, the framework is designed to (i) generate Minutes-style synthetic texts under alternative macro-financial conditions and (ii) evaluate whether such synthetic communications preserve policy-relevant signals, including the ability to infer the direction of policy-rate changes.


\subsection*{Research Question}

This chapter investigates whether large language models can be used as simulation tools for monetary-policy communication and decision support. The main research question is:

\textit{Can a post-trained, reasoning-capable LLM generate FOMC-Minutes-style sections grounded in macroeconomic and financial indicators, and what evidence can the currently archived experiments provide about the resulting text's stylistic fidelity, input grounding, and downstream policy relevance?}

To address this central question, we study the following sub-questions:

\begin{itemize}
\item[1.] To what extent can a fine-tuned LLM reason over macro-financial indicators and extract policy-relevant insights?
\item[2.] Can the framework replicate the structured deliberation and trade-off language observed in FOMC discussions?
\item[3.] How closely do generated sections match the tone, format, and information content of authentic Minutes?
\item[4.] How does model-based decision prediction compare with historical, econometric, and market-implied baselines on held-out FOMC meetings?
\item[5.] Are generated reasoning paths interpretable and consistent with communication patterns in FOMC materials?
\end{itemize}


% Addressing these questions contributes to the literature on AI-enabled macro-financial analysis and central-banking decision support.


\subsection*{Contribution}

This study makes several important contributions to the literature at the intersection of central banking, financial economics, and artificial intelligence:

\begin{itemize}
\item[] \textbf{Methodological Innovation:} We propose a two-step post-training pipeline (SFT + GRPO) to adapt an LLM for data-grounded macro-financial reasoning, together with a downstream alignment procedure for generating structured Minutes-style text.

\item[] \textbf{Checkpointed Research Pipeline:} Rather than treating all post-training outputs as a single ``fine-tuned model,'' we distinguish the analysis checkpoint, the Minutes-style alignment checkpoint, and the decision-specific checkpoint, which allows each empirical result to be mapped to a specific archived model state.

\item[] \textbf{Framework for Synthetic Text Generation:} We propose a framework for generating structured textual data through LLM fine-tuning. This approach translates discrete qualitative information into a continuous distribution, offering new perspectives on risk management strategies in the financial sector.

\item[] \textbf{Extension of LLM Applications:} This study represents one of the first empirical investigations into the use of LLMs as simulation tools for modeling and generating structured textual data. It extends the application of LLMs beyond conventional text analysis and information retrieval tasks.

\item[] \textbf{Advancement of Reasoning LLMs in Decision-Making:} By emphasizing the generation of step-by-step reasoning chains and interpretable decision paths, this study contributes to the growing literature on reasoning-capable LLMs as analytical tools for supporting and simplifying complex decision-making processes.

\item[] \textbf{Practical and Policy-Relevant Implications:} The findings provide a structured benchmark for assessing when LLM-generated policy text is informative enough for simulation use and when the evidence remains exploratory or insufficient for stronger policy claims.

\item[] \textbf{Multi-Dimensional Evaluation:} We evaluate synthetic Minutes not only with text-based similarity metrics, but also with leave-one-out masking tests and downstream economics-based validations, while explicitly documenting that the canonical evidence comes from held-out meeting splits rather than temporal out-of-sample forecasting.
\end{itemize}

This chapter is organized as follows. Section~\ref{ch2:sec:lr} reviews the related literature on synthetic data and large language models. Section~\ref{ch2:sec:2-step-training} outlines the two-step post-training framework. Section~\ref{ch2:sec:the-dataset} describes the dataset construction process and prompt templates. Section~\ref{ch2:sec:base_model} introduces the backbone model, while Sections~\ref{ch2:sec:training} and~\ref{ch2:sec:parameter} describe the fine-tuning method and hyperparameter settings. Section~\ref{ch2:sec:application} presents the evaluation dimensions and empirical tests, and Section~\ref{ch2:sec:result} reports the training and evaluation results. Finally, Section~\ref{ch2:sec:conclusion} concludes the chapter and discusses future research.
